\section{Data and Problem Setup}
\label{sec:data}

\subsection{The ZTBus dataset}

We use the ZTBus dataset introduced by Widmer et
al.~\cite{widmer2023ztbus}, released under CC-BY-4.0. The dataset contains
1,409 trolleybus missions recorded by two HESS lighTram~19 DC articulated
vehicles operated by Verkehrsbetriebe Z\"urich (VBZ) in Zurich,
Switzerland, between April 2019 and December 2022. Each mission is one
complete driving assignment (from pull-out to pull-in) of typically 3--10
hours' duration. Sensor data is recorded at 1~Hz across 26 channels
covering GPS, odometry, electric powertrain, temperature, door state, and
passenger count reported by the on-board Integrated Traffic Control System
(ITCS).

Each mission CSV contains 1~Hz sensor data with timestamps in both ISO and
Unix formats. A separate metadata file lists mission-level attributes:
start/end times, bus number, bus route, and driven distance. The dataset
covers 10 bus routes, with Route~83 accounting for approximately 60\% of
all missions. We use all 1,409 missions in our experiments.

\subsection{The sparse event pattern}

The \texttt{itcs\_numberOfPassengers} channel, the source of our target
variable, is sparse. The ITCS updates passenger count only at stop events,
when doors open and passengers board or alight. At the 1~Hz sampling rate,
only approximately 1\% of rows contain a genuine passenger count; the
remainder are \texttt{NaN}. Stop events occur roughly every 60--120 seconds
of driving time, with some variability depending on route characteristics
(dwell time, stop spacing, time of day).

We forward-fill \texttt{itcs\_numberOfPassengers} within each mission: rows
between stop events inherit the most recent known passenger count. This is
physically justified: passenger count cannot change while the bus is in
motion between stops, and it transforms the 0.9\% usable rows into
approximately 99\% (rows before the first stop of a mission remain
\texttt{NaN} and are dropped). After forward-fill and dropping pre-first-stop
rows, we have 47,628,626 usable rows across the 1,409 missions.

Critically, we preserve the original \texttt{NaN} pattern as a boolean
column \texttt{is\_stop\_event}. A row is a stop event if and only if its
\texttt{itcs\_numberOfPassengers} was non-\texttt{NaN} before forward-fill.
This marker is what enables the stop-level lag feature computation
described in Section~\ref{sec:methods}. Figure~\ref{fig:leakage} illustrates
the forward-fill staircase and the role of \texttt{is\_stop\_event}.

\subsection{Classification task and target construction}

We formulate the problem as 3-class ridership demand classification at
stop-time granularity. For each row in the dense (forward-filled) stream,
we predict whether the current passenger count is \emph{low}, \emph{medium},
or \emph{high}. The class boundaries are tercile thresholds $(q_1, q_2)$
computed from the training partition of \texttt{itcs\_numberOfPassengers}:
\begin{equation}
y(t) = \begin{cases}
  \text{low}    & \text{if } p(t) \le q_1 \\
  \text{medium} & \text{if } q_1 < p(t) \le q_2 \\
  \text{high}   & \text{if } p(t) > q_2
\end{cases}
\end{equation}
where $p(t)$ is the passenger count at row $t$. On the full dataset,
$q_1 \approx 9$ and $q_2 \approx 19$ passengers. The resulting class
distribution is approximately balanced at 37\% low, 32\% medium, 32\% high
(Figure~\ref{fig:passenger-dist}). Importantly, the tercile boundaries are
computed \emph{exclusively} on the training data to prevent test-set
information from leaking into the target definition.

\begin{figure}[t]
  \centering
  \includegraphics[width=\columnwidth]{01_passenger_distribution.png}
  \caption{Distribution of passenger counts across all ZTBus stop events,
    with tercile boundaries $q_1 \approx 9$ and $q_2 \approx 19$ marked.
    The resulting 3-class target is approximately balanced (37/32/32\%).}
  \label{fig:passenger-dist}
\end{figure}

\subsection{Evaluation frameworks}

We evaluate under two splitting strategies:
\begin{itemize}
  \item \textbf{Pooled split.} 80/20 mission-level random split. Train and
    test come from the same distribution. This measures in-distribution
    generalization.
  \item \textbf{Temporal split.} Missions starting before 2022-01-01 are
    training (811 missions, 57.6\%, covering April 2019 through
    December 2021); missions starting in 2022 are test (598 missions,
    42.4\%, covering January through December 2022). The 57/42
    proportion is closer to balanced than a conventional 80/20 temporal
    cut because 2022 was the single most data-rich year in the dataset
    (598 missions, compared to 147 in 2019, 174 in 2020, and 490 in
    2021). This measures generalization to future data and is the more
    realistic deployment scenario.
\end{itemize}

We report macro-averaged F1 score and balanced accuracy under both
frameworks, plus $3\times 3$ confusion matrices and per-class
precision/recall. Macro-F1 is our primary metric because it weights classes
equally, preventing the majority class from dominating the evaluation under
any class imbalance.

\subsection{Mission-level splits prevent row-level leakage}

All splits are at the mission level, not the row level. Since a single
mission generates tens of thousands of 1~Hz rows, a row-level random split
would scatter the same mission across train and test, giving the model
indirect access to other rows from the same trip (highly correlated
neighbors via rolling features, nearby stops, and shared contextual
features). Mission-level splits ensure that a mission is entirely in one
partition.
